\section{Related Work}
\label{sec:related}

\textbf{CI Merge Queues, Batching, and Prioritization.}
Large-scale industrial CI systems introduced speculative execution and dynamic batching to sustain mainline velocity. Uber's SubmitQueue~\cite{ananthanarayanan2019keeping} pioneered speculative merge trees to keep large monorepo mainlines green, while recent industrial studies examine lean merge pipelines~\cite{juloori2025scale} and PR prioritization to improve merge throughput~\cite{jungwirth2025improving}. Elbaum et al.~\cite{elbaum2014techniques} established regression batching techniques at Google, later extended by Bavand and Rigby~\cite{bavand2021mining} for failure-aware dynamic batching and Fallahzadeh et al.~\cite{fallahzadeh2023accelerating} for parallel batch testing with bisection. Complementary CI acceleration works select subsets of builds or tests via transition analysis~\cite{leong2019assessing}, reinforcement learning~\cite{spieker2017reinforcement}, hybrid build/test prediction~\cite{jin2023hybridcisave}, or pipeline-aware selection~\cite{schwendner2025practical}. Unlike these approaches, SSMQ operates at the \emph{PR scheduling and batch composition} layer using dependency-graph disjointness and causal blame attribution, making it orthogonal and composable with regression test selection.

\textbf{Flaky Tests and Noisy Culprit Finding.}
Empirical studies demonstrate that flaky tests are pervasive across languages and UI suites~\cite{luo2014empirical,lam2020large,romano2021empirical,tahir2023test} and often co-occur systemically~\cite{parry2025systemic}. Detection frameworks such as DeFlaker~\cite{bell2018laker} and iDFlakies~\cite{lam2019idflakies} identify flaky tests via differential coverage and order perturbation. For culprit localization, Ziftci and Reardon~\cite{ziftci2017broke} automated regression attribution at Google scale, and Henderson et al.~\cite{henderson2023flake} formulated flake-aware Bayesian noisy bisection. SSMQ adapts probabilistic flakiness modeling into pre-merge speculative batch control and empirically characterizes where selective reruns outperform or lag immediate fault isolation.

\textbf{Semantic Merge Conflicts, Change Impact Analysis, and LLM Fault Localization.}
Prior work has explored predicting merge conflicts~\cite{owhadikareshk2019predicting}, using PR-level change impact analysis to enhance code review~\cite{gmen2025enhanced}, and resolving textual or semantic conflicts with neural transformers and pre-trained language models~\cite{zhang2022trained,svyatkovskiy2022program}. Recent studies also demonstrate LLM-driven fault localization from crash and test stack traces~\cite{jambigi2025fault}. SSMQ bridges these traditionally separate communities by using change-impact graphs proactively to prevent semantic conflicts inside merge batches and coupling lightweight causal blame directly into merge-queue recovery.
